% !TeX program = xelatex
% =============================================================================
%  黎黄江 (GIANG LE HUYNH) — 中文简历 / Chinese CV
% -----------------------------------------------------------------------------
%  ⚠ BẮT BUỘC compile bằng XeLaTeX. pdfLaTeX KHÔNG render được chữ Hán.
%
%    Overleaf : bấm "Menu" (góc trên trái) -> phần Settings -> Compiler
%               -> chọn XeLaTeX -> đóng menu -> Recompile
%               (Overleaf KHÔNG tự đọc dòng "% !TeX program" ở trên;
%                bắt buộc phải đổi ở dropdown Compiler)
%    Máy local: xelatex cv_china.tex
%
%  Font: fandol — có sẵn trong TeX Live / Overleaf, không cần cài thêm.
%  Nếu log báo "fontset 'fandol' is unavailable in current mode"
%    => vẫn đang dùng pdfLaTeX, chưa đổi compiler.
% =============================================================================

\documentclass[10pt, letterpaper]{article}

% --- 中文支持 (必须放在最前) ---
\usepackage[UTF8, fontset=fandol]{ctex}

\usepackage{tabularx}
\usepackage{booktabs}
\usepackage{etoolbox}
\usepackage[
    ignoreheadfoot,
    top=2 cm,
    bottom=2 cm,
    left=2 cm,
    right=2 cm,
    footskip=1.0 cm,
]{geometry}
\usepackage[explicit]{titlesec}
\usepackage{array}
\usepackage[dvipsnames]{xcolor}
\definecolor{primaryColor}{RGB}{0,0,0}
\usepackage{enumitem}
\usepackage{fontawesome5}
\usepackage{amsmath}
\usepackage[
    pdftitle={黎黄江 - AI 工程师 简历},
    pdfauthor={Giang Le Huynh},
    pdfcreator={LaTeX},
    colorlinks=true,
    urlcolor=primaryColor
]{hyperref}
\usepackage[pscoord]{eso-pic}
\usepackage{calc}
\usepackage{bookmark}
\usepackage{lastpage}
\usepackage{changepage}
\usepackage{paracol}
\usepackage{ifthen}
\usepackage{needspace}

% --- 设置 ---
\AtBeginEnvironment{adjustwidth}{\partopsep0pt}
\pagestyle{empty}
\setcounter{secnumdepth}{0}
\setlength{\parindent}{0pt}
\setlength{\topskip}{0pt}
\setlength{\columnsep}{0.15cm}

\makeatletter
\let\ps@customFooterStyle\ps@plain
\patchcmd{\ps@customFooterStyle}{\thepage}{
    \color{gray}\textit{\small 黎黄江 (Giang Le Huynh) - 第 \thepage{} 页 / 共 \pageref*{LastPage} 页}
}{}{}
\makeatother
\pagestyle{customFooterStyle}

\titleformat{\section}{
    \needspace{4\baselineskip}
    \Large\color{primaryColor}
}{
}{
}{
    \textbf{#1}\hspace{0.15cm}\titlerule[0.8pt]\hspace{-0.1cm}
}[]

\titlespacing{\section}{-1pt}{0.1 cm}{0.1 cm}

\newenvironment{highlights}{
    \begin{itemize}[
        topsep=0.10 cm,
        parsep=0.10 cm,
        partopsep=0pt,
        itemsep=0pt,
        leftmargin=0.1 cm + 10pt
    ]
}{
    \end{itemize}
}

\newenvironment{onecolentry}{
    \begin{adjustwidth}{0.2 cm + 0.00001 cm}{0.2 cm + 0.00001 cm}
}{
    \end{adjustwidth}
}

\newenvironment{threecolentry}[3][]{
    \onecolentry
    \def\thirdColumn{#3}
    \setcolumnwidth{1.6 cm, \fill, 4.5 cm}
    \begin{paracol}{3}
    {\raggedright #2} \switchcolumn
}{
    \switchcolumn \raggedleft \thirdColumn
    \end{paracol}
    \endonecolentry
}

\newenvironment{header}{
    \setlength{\topsep}{0pt}\par\kern\topsep\centering\color{primaryColor}\linespread{1.5}
}{
    \par\kern\topsep
}

\newcommand{\placelastupdatedtext}{%
  \AddToShipoutPictureFG*{%
    \put(
        \LenToUnit{\paperwidth-2 cm-0.2 cm+0.05cm},
        \LenToUnit{\paperheight-1.0 cm}
    ){\vtop{{\null}\makebox[0pt][c]{
        \small\color{gray}\textit{更新于 2026 年 7 月}\hspace{\widthof{更新于 2026 年 7 月}}
    }}}%
  }%
}%

\let\hrefWithoutArrow\href
\renewcommand{\href}[2]{\hrefWithoutArrow{#1}{\ifthenelse{\equal{#2}{}}{ }{#2 }\raisebox{.15ex}{\footnotesize \faExternalLink*}}}

\begin{document}
    \newcommand{\AND}{\unskip
        \cleaders\copy\ANDbox\hskip\wd\ANDbox
        \ignorespaces
    }
    \newsavebox\ANDbox
    \sbox\ANDbox{}

    \placelastupdatedtext
    \begin{header}
        \fontsize{30 pt}{30 pt}
        \textbf{黎黄江}\ \ \fontsize{16 pt}{16 pt}\textbf{(GIANG LE HUYNH)}
        \\
        \vspace{0.1 cm}

        \fontsize{18 pt}{18 pt}
        \textbf{AI 工程师 / 提示词工程师}

        \vspace{0.1 cm}

        \kern 0.25 cm%
        \mbox{\hrefWithoutArrow{mailto:lhgiang040504@gmail.com}{{\footnotesize\faEnvelope[regular]}\hspace*{0.13cm}lhgiang040504@gmail.com}}%
        \kern 0.25 cm%
        \AND%
        \kern 0.25 cm%
        \mbox{\hrefWithoutArrow{tel:+84812206773}{{\footnotesize\faPhone*}\hspace*{0.13cm}+84 812 206 773 (WhatsApp)}}%
        \kern 0.25 cm%
        \AND%
        \kern 0.25 cm%
        \mbox{\hrefWithoutArrow{https://linkedin.com/in/lhgiang040504}{{\footnotesize\faLinkedinIn}\hspace*{0.13cm}lhgiang040504}}%
        \kern 0.25 cm%
        \AND%
        \kern 0.25 cm%
        \mbox{\hrefWithoutArrow{https://github.com/lhgiang040504}{{\footnotesize\faGithub}\hspace*{0.13cm}lhgiang040504}}%
        \\
        \vspace{0.05 cm}
        \fontsize{9 pt}{9 pt}
        \mbox{\hrefWithoutArrow{https://lhgiang040504.github.io/lehuynhgiang/?lang=zh}{作品集（中文版）: lhgiang040504.github.io/lehuynhgiang/?lang=zh}}%
    \end{header}

    \vspace{0.1 cm}

% ------------------------------------------------------------------------------------------

\section{个人简介}
\begin{onecolentry}
    计算机科学专业应届生（人工智能与信息安全方向），现任 AI 工程师，负责为企业级 B2B 客户（新加坡、日本）交付生产环境的 AI 智能体产品，覆盖需求梳理、架构设计、全栈开发、部署到线上运维的完整流程。

    \vspace{0.15 cm}

    核心专长为\textbf{提示词架构、大模型评测与安全护栏（Guardrails）}：已发表 2 篇 IEEE 论文（其中 1 篇为第一作者），研究方向即为\textbf{结构化提示词流程设计与降低大模型幻觉}。既做研究，也能把系统真正上线并持续运行。

    \vspace{0.15 cm}

    \textbf{语言说明（如实告知）}：英语可胜任专业沟通（B2，日常与新加坡及日本客户对接）；越南语为母语；\textbf{中文目前为零基础，已开始学习，目标 6 个月内达到 HSK 3}。此前曾在无日语基础的情况下，独立完成日语\textbf{敬语}语音客服系统的提示词设计与结构，由母语同事复核语气，本人负责逻辑、约束与测试体系。
\end{onecolentry}

% ------------------------------------------------------------------------------------------

\section{工作经历}
\begin{onecolentry}
    \textbf{AI 工程师 — FunAI (Functional AI Partners)} \hfill \textit{2026 年 5 月 -- 至今}
    \begin{highlights}
\item 主导为新加坡、日本企业客户交付\textbf{多个生产环境「数字员工」AI 智能体}，从与客户共同梳理需求，到架构设计、开发、部署与运维全流程负责
\item 基于 \textbf{LangGraph + FastAPI} 设计 WhatsApp 对话式智能体，自动化多步骤客户入驻流程；底层采用 Postgres 工作队列 + \textbf{分布式锁}，并支持人工介入（human-in-the-loop）升级，在 \textbf{约 30 路并发会话}下保持稳定
\item 交付自主文档分类智能体，通过 \textbf{16+ 个 n8n 编排工作流}完成业务文档的接入、分类与归档，并配套基于安全参数化 SQL 的自然语言检索助手
\item 构建\textbf{多智能体日语语音 AI} 客服系统（ElevenLabs、Firebase、Twilio）：负责敬语提示词设计、KYC 身份核验，以及守护每次发版的 \textbf{68 条回归测试用例}
\item 开发分布式\textbf{人脸检测与聚类}流水线（PyTorch、YOLO、InsightFace、HDBSCAN），含 GPU 任务编排，并封装为可处理 \textbf{5 GB+} 压缩包上传的全栈平台
\end{highlights}
\end{onecolentry}

\vspace{0.2 cm}

\begin{onecolentry}
    \textbf{AI 研究员 — InsecLab（校级实验室）} \hfill \textit{2025 年 5 月 -- 2026 年 3 月}
    \begin{highlights}
\item 发表 \textbf{2 篇 IEEE 论文}（1 篇第一作者），主题为基于大语言模型与多模态方法的智能合约漏洞检测
\item 设计\textbf{六阶段结构化提示词流程}（角色定义 → 验证 → 批判 → 推理链），提升推理一致性并降低大模型审计中的幻觉
\item 在统一提示词下横向评测多个大模型（OpenAI o 系列、Llama 等），并搭建可扩展的向量嵌入与图数据流水线，支撑基准测试与消融实验
\end{highlights}
\end{onecolentry}

\vspace{0.2 cm}

\begin{onecolentry}
    \textbf{AI 研究合作 — Pi Labs AI（印度公司，远程）} \hfill \textit{2025 年 9 月 -- 2025 年 12 月}
    \begin{highlights}
\item 在产业合作项目中构建电信数据（CDR/IPDR）异常检测方案，并参与定义问题范围与评估标准
\item 评测 RAG 框架（KAG），探索基于图结构及多模态（CDR + 语音/文本）的欺诈检测方案
\end{highlights}
\end{onecolentry}

% ------------------------------------------------------------------------------------------

\section{学术论文}
\begin{onecolentry}
    \href{https://ieeexplore.ieee.org/abstract/document/11133797}{基于提示词的智能合约分析 (Prompt-Based Smart Contract Analysis)} \hfill \textit{共同作者，IEEE MAPR 2025}
    \begin{highlights}
\item 设计并实现\textbf{六阶段结构化提示词流程}，提升推理一致性、降低幻觉；横向评测 OpenAI o 系列与 Llama 等多个模型
    \end{highlights}
\end{onecolentry}

\vspace{0.15 cm}

\begin{onecolentry}
    \href{https://ieeexplore.ieee.org/abstract/document/11231648}{面向智能合约安全的多模态方法 (Multimodal Approach for Smart Contract Security)} \hfill \textit{\textbf{第一作者}，IEEE ISCIT 2025}
    \begin{highlights}
\item 多模态深度学习流水线，漏洞检测 \textbf{F1 = 0.823}、\textbf{准确率 96.73\%}；基准测试五种特征融合策略
    \end{highlights}
\end{onecolentry}

\vspace{0.15 cm}

\begin{onecolentry}
    融合大语言模型与异构图的智能合约安全研究 \hfill \textit{\textbf{第一作者}，已投稿 2026}
\end{onecolentry}

% ------------------------------------------------------------------------------------------

\section{项目经历}
\begin{onecolentry}
    \textbf{模块化检索增强生成（RAG）系统} \hfill \textit{2025 年 1 月 -- 2025 年 5 月} \\
    \href{https://github.com/lhgiang040504/sematic_search}{GitHub}

    \textbf{技术栈}：Python、FastAPI、MongoDB、LangChain、SentenceTransformers、Chroma

    \begin{highlights}
\item 设计支持\textbf{多种查询策略}的模块化 RAG 流水线（查询分解 + RRF 倒数排序融合）
\item 在 \textbf{MS MARCO}、\textbf{Cranfield} 基准数据集上，相较基线检索质量提升 \textbf{22\%}，P95 延迟 < 1.8 秒
    \end{highlights}
\end{onecolentry}

\vspace{0.2 cm}

\begin{onecolentry}
    \textbf{简历—岗位智能匹配系统} \hfill \textit{2024 年 11 月 -- 2025 年 1 月} \\
    \href{https://github.com/lhgiang040504/Job-Resume-Matching-Project}{GitHub}

    \textbf{技术栈}：Python、FastAPI、MongoDB、React、LangChain、LLaMA 3、SentenceTransformers

    \begin{highlights}
\item 全栈 RAG 系统，自研\textbf{简历—岗位相似度打分机制}，上下文 precision@5 达 \textbf{0.87}
\item 基于 400+ 条爬取岗位描述与真实简历数据，推荐响应延迟 < 5 秒，并接入延迟与 token 用量监控
    \end{highlights}
\end{onecolentry}

% ------------------------------------------------------------------------------------------

\section{教育背景}
\begin{threecolentry}{\textbf{本科}}{2022 -- 2026}
        \textbf{胡志明市国家大学 — 信息技术大学 (UIT, VNU-HCM)} \\
        \textbf{专业}：计算机科学 \\
        \textbf{方向}：人工智能与信息安全 \\
        \textbf{相关课程}：人工智能、机器学习、深度学习、数据挖掘、MLOps、AI 系统
\end{threecolentry}

% ------------------------------------------------------------------------------------------

\section{技术能力}
\begin{onecolentry}
\textbf{提示词与大模型工程：} 提示词架构设计、系统提示词（System Prompt）、思维链 CoT / ReAct / Few-shot、\textbf{结构化输出}（JSON / XML / Markdown 约束）、工具与函数调用、提示词缓存、\textbf{大模型评测}（测试集构建、指标设计、多模型对比）、\textbf{安全护栏}（防越狱、防提示注入、降低幻觉）

\vspace{0.1 cm}

\textbf{大模型与智能体：} 多智能体与智能体工作流（LangGraph、ReAct）、RAG / 高级检索、语音 AI（ElevenLabs）、Anthropic Claude、OpenAI、Google Gemini、LLaMA

\vspace{0.1 cm}

\textbf{AI 开发工具：} Claude Code、Anthropic Console / Workbench、OpenAI Playground、Codex、Cursor、MCP（Model Context Protocol）服务配置、n8n

\vspace{0.1 cm}

\textbf{全栈与云：} Next.js / React、TypeScript、FastAPI、Supabase / pgvector、Vercel、Railway、Firebase / GCP、Cloudflare R2

\vspace{0.1 cm}

\textbf{MLOps 与生产化：} Docker、CI/CD、SSE 流式输出、监控与可观测性（延迟、成本）、实验追踪

\vspace{0.1 cm}

\textbf{检索与数据：} 混合检索、向量检索、RRF、BM25、SentenceTransformers、SQL、MongoDB

\vspace{0.1 cm}

\textbf{语言：} 英语（B2 中高级，可胜任专业沟通）· 越南语（母语）· 中文（零基础，正在学习，目标 6 个月内 HSK 3）
\end{onecolentry}

\end{document}
